\documentclass[a4paper,11pt,dvipsnames]{article}
\usepackage[a4paper,vmargin={15mm,20mm},hmargin={20mm,15mm}, includefoot, bottom=3.5cm]{geometry}

\setlength\parindent{0pt}

\usepackage{amsmath,amsthm,amssymb,amsfonts}
\usepackage[utf8]{inputenc}
\usepackage[T1]{fontenc}
\usepackage[spanish]{babel}

\usepackage[per-mode=fraction, separate-uncertainty=true]{siunitx}
\usepackage{physics}
\AtBeginDocument{\RenewCommandCopy\qty\SI}

\usepackage{pgfplots}
\usepackage{tikz}
\tikzset{axis/.style={thick,-latex}}
\usetikzlibrary{decorations}
\usetikzlibrary{decorations.pathmorphing}
\usetikzlibrary{shapes.geometric}
\usetikzlibrary{arrows}
\usetikzlibrary{arrows.meta}
\usetikzlibrary{calc}
\pgfplotsset{compat=newest}
\usetikzlibrary{patterns}

\tikzstyle{axisarrow} = [-{Latex[inset=0pt,length=5pt]}]

\usepackage{enumerate}
\usepackage{fancybox}
\usepackage{fancyhdr}
\setlength{\headheight}{60pt}
\pagestyle{fancy}
\fancyhead{}
\rhead{convocatoria \Large 2026}
\lhead{\Large FCR}
\chead{\Large Chatbot IA Hub}

\begin{document}
\thispagestyle{fancy}

\noindent
\begin{center}
{\LARGE \vspace*{0.5cm} Ejercicios propuestos a los estudiantes}\par\medskip \textit{(Material para docentes)}
\end{center}

\vspace{20pt}

\begin{itemize}
 
\item \textbf{calorimetria-001 :} Aceite y hielo, buscar masa del aceite. \par\medskip

\begin{center}
\begin{minipage}{0.60\linewidth}    
\textbf{Ejercicio:} Un calorímetro ideal y adiabático contiene, inicialmente, una masa de aceite de oliva a $80^\circ\text{C}$. Se introduce en el calorímetro un bloque de $100\text{ g}$ de hielo a $0^\circ\text{C}$. Luego de un rato, se encuentra que el hielo se ha fundido por completo, y tanto el agua como el aceite se encuentran a $40^\circ\text{C}$. Sabiendo que el calor específico del aceite de oliva es $c_\text{aceite} = 0.4\textrm{\,cal/(g $^\circ$C)}$, determinar la masa de aceite de oliva contenida en el calorímetro.
\end{minipage}
\end{center}


\item \textbf{termodinamica-001 :} Trabajo de un ciclo con proceso desconocido. \par\medskip

\begin{center}
\begin{minipage}{0.60\linewidth}
  \textbf{Ejercicio:} La figura muestra un ciclo termodinámico para un gas ideal monoatómico. Se sabe que la temperatura del gas en el estado $a$ es $260\text{ K}$, y que el gas cede $800\text{ J}$ de calor en el proceso $c\to a$. Calcular el trabajo neto realizado por el gas en el ciclo.
  
  \begin{center}
  \begin{tikzpicture}[scale=1]
    \begin{axis}[
      axis x line=bottom,
      axis y line=left,
      xmin=1.5, xmax=3.5,           %min y max para los ejes, NO PARA EL DOMINIO
      ymin=0.25, ymax=2.5, 
      xlabel={volumen [$\si{\litre}$]},
      ylabel={presión [atm]},
      xtick={2,3},
      xticklabels={$2$,$3$},
      ytick={0.75,2},
      yticklabels={$0.75$,$2$}
      ]
    % \draw [color=blue, very thick][-latex](160pt,20pt)..controls(100pt,25pt) and (60pt,45pt)..(42pt,58pt);
    \draw [color=blue, very thick][-latex](2,2)--(2.5,2);
    \draw [color=blue, very thick](2.5,2)--(3,2);
    \draw [color=blue, very thick][-latex] (2,0.75)--(2,1.5);
    \draw [color=blue, very thick] (2,1.5)--(2,2);

    % Parábola de c (3,2) a a (2,0.75): P(V) = 0.25*V^2 - 0.25
    \addplot [color=blue, very thick, latex-, samples=50, domain=2.5:3] {1.25*(x-2)^2 + 0.75};
    \addplot [color=blue, very thick, samples=50, domain=2:2.5] {1.25*(x-2)^2 + 0.75};

    \draw [dashed] (2,0)--(2,0.75);
    \draw [dashed] (3,0)--(3,2);
    \draw [dashed] (0,0.75)--(2,0.75);
    \draw [dashed] (0,2)--(2,2);
    \filldraw(2,0.75)circle(2pt) (2,2)circle (2pt)(3,2)circle(2pt);
    \draw (2,0.75) node [below left] {$a$};
    \draw (2,2) node [above left] {$b$};
    \draw (3,2) node [above right] {$c$};
    \end{axis}
  \end{tikzpicture}
  \end{center}

\end{minipage}
\end{center}

\end{itemize}
  
\end{document}